\par\noindent\SourceLeaf{49}%
\hspace*{\parindent}Terminus generalis $r^{\text{tus}}$
\[
= \frac{1}{2^{2m-1}}\,{}^{2m}\mathfrak{B}^{(m+r)}\cos 2r\varphi\cos\varphi,
\]
\par\noindent in partes diremptus, dabit
\[
\frac{1}{2^{2m-1}}\,{}^{2m}\mathfrak{B}^{(m+r)} . \frac{1}{2}\cos(2r-1)\varphi
+ \frac{1}{2^{2m-1}}\,{}^{2m}\mathfrak{B}^{(m+r)} . \frac{1}{2}\cos(2r+1)\varphi.
\]

\par\noindent Igitur sola multipla imparia hic occurrunt, quorum bina contigua ex quolibet
datae seriei termino generantur. Si igitur velim novae seriei terminum
$(r-1)^{\text{tum}}$, quem in $\cos(2r-1)\varphi$ ductum fore adparet, ex
evolutae seriei termino $(r-1)^{\text{to}}$ partem posteriorem
\[
= \frac{1}{2^{2m}}\,{}^{2m}\mathfrak{B}^{(m+r-1)}\cos(2r-1)\varphi,
\]
ex $r^{\text{to}}$ autem partem priorem
\[
\frac{1}{2^{2m}}\,{}^{2m}\mathfrak{B}^{(m+r)}\cos(2r-1)\varphi
\]
iungendo, obtinerem:
\[
\frac{1}{2^{2m}}\,{}^{2m+1}\mathfrak{B}^{(m+r)}\cos(2r-1)\varphi.
\]

\par\noindent Sic evadet ipsius $\cos\varphi^{2m+1}$ terminus generalis $(r-1)^{\text{tus}}$
qualis hic propositus est, aut, si malis, pro $r$ substituendo $r+1$, terminus
$r^{\text{tus}}$
\[
\frac{1}{2^{2m}}\,{}^{2m+1}\mathfrak{B}^{(m+1+r)}\cos(2r+1)\varphi.
\]

\par\noindent Atque iam adparet hanc legem, ipsius $\cos\varphi^{2m+1}$ terminum
$r^{\text{tum}}$ sistentem, cum adsumta illa pro $\cos\varphi^{2m-1}$ termino
$r^{\text{to}}$
\[
= \frac{1}{2^{2m-2}}\,{}^{2m-1}\mathfrak{B}^{(m+r)}\cos(2r+1)\varphi,
\]
ita congruere, ut, simulac in hac loco $m$ substituatur $m+1$, illa sponte
prodeat.

\par\noindent\hspace*{\parindent}Esse igitur legem tam pro paribus quam pro imparibus ipsius $\cos\varphi$
potestatibus generalem, hoc progressu indefinito demonstratione munito, per se
patet.

5. Quamvis ad disquisitionem propositam haud pertineat, liceat hoc loco
monere, ex lege pro cosinuum potestatibus per dati anguli multiplos exhibitis
similem pro sinuum potestatibus statim exsurgere, idque sola observatione esse
$\cos(90^0-\varphi) = \sin\varphi$.
